\documentclass[11pt]{article}
\usepackage{amsmath,amssymb,booktabs,geometry,graphicx,hyperref,url}
\usepackage{subcaption,placeins}
\hypersetup{hidelinks}
\geometry{margin=1in}
\setlength{\parskip}{0.45em}
\setlength{\parindent}{0pt}
\title{Stage 1: Mathematical Formulation Notes}
\date{\today}
\begin{document}
\maketitle

\section{Stage 1: Mathematical Formulation Notes}
\subsection{Purpose}
This stage establishes the mathematical and notational contract used by every
later solver, diagnostic, and surrogate experiment. An American option may be
exercised at any time up to maturity, so its value must never fall below the
immediate-exercise payoff. That additional decision makes the problem an
obstacle problem with an unknown early-exercise boundary, rather than the
fixed-boundary valuation problem associated with a European option.

\subsection{Method}
Time is expressed as time to maturity, $\tau=T-t$, so the computation starts
from the payoff at $\tau=0$ and marches forward to the valuation horizon.
Numerical experiments use $K=1$ unless otherwise stated, while reports retain
normalized quantities such as $S/K$, $U/K$, and $B/K$ so the interpretation
does not depend on this implementation choice.

\subsection{Mathematical Formulation}
\[
  dS_t=(r-q)S_t\,dt+\sigma S_t\,dW_t,
\qquad
  \mathcal{L}U=
  \frac{1}{2}\sigma^2S^2U_{SS}+(r-q)SU_S-rU.
\]
Here \(r\) is the risk-free rate, \(q\) is the continuous dividend yield,
\(\sigma\) is volatility, and \(U(S,\tau)\) is the American option value. For
payoff \(\Phi(S)\), the time-to-maturity variational inequality is
\[
  U\geq \Phi,\qquad
  U_\tau-\mathcal{L}U\geq0,\qquad
  (U-\Phi)(U_\tau-\mathcal{L}U)=0.
\]
The first inequality is the payoff obstacle. In the exercise region
\(U=\Phi\); in the continuation region \(U>\Phi\) and the Black--Scholes PDE
holds. Their interface is the free boundary \(B(\tau)\). The continuation
premium
\[
  C(S,\tau)=U(S,\tau)-\Phi(S)\geq0
\]
measures the value of retaining the option rather than exercising immediately.

\subsection{Implementation Details}
The solver operates on a spot-time grid \((S_i,\tau_n)\). Later surrogate
datasets use
\[
  x=\bigl(\log(S/K),\tau/T,r,q,\sigma,T,\mathbf{1}_{\mathrm{call}}\bigr),
\]
with labels reported as \(U/K\), \(C/K\), \(\Delta\), and \(K\Gamma\). This
stage defines those quantities only; it does not yet construct data or train a
model.

\subsection{How to Read the Matrix Calculations}
The finite-difference implementation converts the continuous operator into a
matrix acting on interior grid values. A three-point stencil approximates the
first and second spot derivatives, producing the lower, diagonal, and upper
coefficients of a tridiagonal operator \(L_h\). Boundary nodes are prescribed
separately and contribute only to the first and last interior right-hand-side
entries. The key reading path is
\[
  \text{PDE operator}
  \longrightarrow \text{three-point stencil}
  \longrightarrow L_h
  \longrightarrow
  \left\{
  \begin{aligned}
    A&=I-\frac{\Delta\tau}{2}L_h,\\
    B&=I+\frac{\Delta\tau}{2}L_h.
  \end{aligned}
  \right.
\]
The matrices \(A\) and \(B\) are the new- and old-time-level components of the
Crank--Nicolson step. The American obstacle then turns the linear update into a
linear complementarity problem, which PSOR handles by alternating relaxation
with projection above \(\Phi\). Translating this formulation into reproducible
arrays, diagnostics, figures, and LaTeX explanations is a central
scientific-computing contribution of the project.

\subsection{Validation and Diagnostics}
Before implementation, the notation audit checks the time direction, the sign
of \(U_\tau-\mathcal{L}U\), put and call payoff conventions, boundary
orientation, and units of \(\Delta\) and \(\Gamma\). These checks prevent
sign-consistent code from solving the wrong mathematical problem.

\subsection{Results}
The output of this stage is an implementation-ready formulation with consistent
normalization, matrix notation, and obstacle terminology. It produces no
numerical accuracy or surrogate-performance evidence.

% Auto-generated by experiments/19_report_visualization_enrichment.py
\section*{Visual Evidence Gallery}
\begin{figure}[!htbp]
\centering
\begin{subfigure}[t]{0.48\linewidth}
\centering
\includegraphics[width=\linewidth]{results/figures/report_enrichment/stage01_math_notes_payoff_curves.pdf}
\caption{Put and call payoff curves establish the obstacle that all American option values must respect.}
\end{subfigure}
\hfill
\begin{subfigure}[t]{0.48\linewidth}
\centering
\includegraphics[width=\linewidth]{results/figures/report_enrichment/stage01_math_notes_exercise_continuation_regions.pdf}
\caption{The exercise and continuation regions show why the early-exercise boundary is a first-class diagnostic rather than a decorative price threshold.}
\end{subfigure}
\end{figure}

\begin{figure}[!htbp]
\centering
\begin{subfigure}[t]{0.48\linewidth}
\centering
\includegraphics[width=\linewidth]{results/figures/report_enrichment/stage01_math_notes_complementarity_geometry.pdf}
\caption{The variational inequality couples obstacle contact with zero PDE residual in continuation.}
\end{subfigure}
\hfill
\begin{subfigure}[t]{0.48\linewidth}
\centering
\includegraphics[width=\linewidth]{results/figures/report_enrichment/stage01_math_notes_continuation_premium.pdf}
\caption{The continuation premium C equals U minus Phi and separates immediate-exercise payoff from the value of waiting; this motivates the later price-only surrogate.}
\end{subfigure}
\end{figure}

\begin{figure}[!htbp]
\centering
\begin{subfigure}[t]{0.48\linewidth}
\centering
\includegraphics[width=\linewidth]{results/figures/report_enrichment/stage01_math_notes_spot_grid_stencil.pdf}
\caption{Three neighboring spot nodes show how first and second derivatives are approximated on the grid.}
\end{subfigure}
\hfill
\begin{subfigure}[t]{0.48\linewidth}
\centering
\includegraphics[width=\linewidth]{results/figures/report_enrichment/stage01_math_notes_coefficient_decomposition.pdf}
\caption{The lower, diagonal, and upper coefficients show how diffusion, drift, and discounting enter one row of the finite-difference operator.}
\end{subfigure}
\end{figure}

\begin{figure}[!htbp]
\centering
\begin{subfigure}[t]{0.48\linewidth}
\centering
\includegraphics[width=\linewidth]{results/figures/report_enrichment/stage01_math_notes_interior_vector_layout.pdf}
\caption{The solver matrix acts on interior unknowns while boundary values enter through the right hand side.}
\end{subfigure}
\hfill
\begin{subfigure}[t]{0.48\linewidth}
\centering
\includegraphics[width=\linewidth]{results/figures/report_enrichment/stage01_math_notes_annotated_operator_matrix.pdf}
\caption{The tridiagonal matrix layout maps the stencil coefficients into rows of the discrete operator.}
\end{subfigure}
\end{figure}

\begin{figure}[!htbp]
\centering
\begin{subfigure}[t]{0.48\linewidth}
\centering
\includegraphics[width=\linewidth]{results/figures/report_enrichment/stage01_math_notes_cn_matrix_pair.pdf}
\caption{The Crank--Nicolson matrices A and B split the same spatial operator across the new and old time levels; the diagram explains construction rather than adding accuracy evidence.}
\end{subfigure}
\hfill
\begin{subfigure}[t]{0.48\linewidth}
\centering
\includegraphics[width=\linewidth]{results/figures/report_enrichment/stage01_math_notes_lcp_vector_form.pdf}
\caption{The vector inequalities connect the payoff obstacle with the matrix residual at every interior node.}
\end{subfigure}
\end{figure}

\begin{figure}[!htbp]
\centering
\begin{subfigure}[t]{0.48\linewidth}
\centering
\includegraphics[width=\linewidth]{results/figures/report_enrichment/stage01_math_notes_time_marching_grid.pdf}
\caption{The finite difference method starts from the payoff and advances through time to maturity.}
\end{subfigure}
\hfill
\begin{subfigure}[t]{0.48\linewidth}
\centering
\includegraphics[width=\linewidth]{results/figures/report_enrichment/stage01_math_notes_rhs_boundary_adjustment.pdf}
\caption{Boundary values affect only the edge equations of the interior linear system.}
\end{subfigure}
\end{figure}

\FloatBarrier


\subsection{Failure Modes and Limitations}
The principal risks are mixing calendar time with time to maturity, reversing
the complementarity sign, applying the put boundary orientation to a
dividend-paying call, or confusing \(C\) with the option price. These risks are
addressed in Stage 2 through analytic checks and explicit matrix diagnostics.

\subsection{Conclusion and Next Step}
The formulation is sufficiently explicit to implement the CN/PSOR benchmark.
The next stage must verify the discrete operator, matrix construction, boundary
terms, and obstacle projection before any risk surface is treated as evidence.

\end{document}
